% Reciprocal full-datum proofs; previously released fragments are untouched.
\annexitem{annex:p10.lawvere-closed-base.partials}{thm:p10-lawvere-closed-base-partials}{Effective partial elements, finite limits and exact denominators}
\begin{proof}
All constructions below use the already established power objects,
internal quantifiers, stable images and effective quotients in the two
original topoi. In particular no desired base functor is a premise.
The general image/quantifier and quotient dependencies are
Theorems~\ref{thm:p10-lawvere-sheaves-predicate-adjoints} and
\ref{thm:p10-lawvere-sheaves-effective-quotients}, respectively.
The universal membership over \(\partial X\) has at most one
member in each fiber, so its projection \(D_X\to\partial X\)
is monic. Its domain classifier is \(d_X\). Intersect two such
universal members over \(\partial X\times\partial X\); the
common-member projection is again monic and has classifier \(e_X\).
The diagonal, symmetry and transitivity equations in the body follow
from intersection of subsingletons: two common members linked by a
third partial member must be equal. These are equations and inequalities
of actual maps to \(T\), at arbitrary parameter objects.

Forward image under \(f:X\to Y\) sends a subsingleton predicate
to a subsingleton predicate, giving \(\partial f\). Its domain
is unchanged: the partial graph maps isomorphically to its image over
the same monic domain. Consequently
\[
 d_Y\partial f=d_X,\qquad
 e_X(p,q)\le e_Y(\partial f(p),\partial f(q)).             \tag{1}
\]
Images of composites equal composite images, by image minimality and
the epi/mono factorization, so \(\partial\) is a functor.
If \(f\) is monic, the inequality in (1) is equality.

The finite-product identifications of \(u_*\) and meet preservation
of \(h\) carry all these equations and inequalities to the hatted
maps. We do not assume that \(u_*\) preserves images or that its
image of a mono is classified by the hatted predicate. On \(A_X\)
the diagonal of \(\widehat e_X\) is true; symmetry and transitivity
hold. Thus \(R_X\) is an equivalence relation and its effective
quotient \(q_X:A_X\to H_hX\) has exactly \(R_X\) as kernel
pair. Equation (1) gives \(A_f:A_X\to A_Y\) preserving that
kernel pair. The coequalizer universal property yields the unique
\(H_hf\) with \(H_hf\,q_X=q_YA_f\). Epi cancellation gives
identity and composition laws. The same property is the claimed
natural bijection between maps out of the quotient and maps equal
on \(R_X\); it is valid for every target \(Z\).

Singletons define \(s_X:X\to\partial X\) with domain true,
so \(\eta_X=q_Xu_*s_X\) is defined. Image functoriality makes
it natural. Its equality relation is exactly the newly defined
predicate \(h\,u_*\chi_{\Delta_X}\): this follows from
\(e_X(s_Xx,s_Xx')=\chi_{\Delta_X}(x,x')\) and effectiveness.
In addition there is an actual truth map \(\tau:H_hT\to T_0\)
extending \(h\). On \(A_T\) it is
\(p\mapsto\widehat e_T(p,u_*(s_T\mathsf{true}))\),
using \(u_*1\cong1_0\).
If \(R_T(p,q)\) holds, transitivity in both directions makes
these two values equal. It therefore descends uniquely; and
\(e_T(s_Tt,s_T\mathsf{true})=t\), since the classifier of
\(\mathsf{true}:1\to T\) is the identity of \(T\).
Thus \(\tau\eta_T=h\). No claim that \(\tau\) is an
isomorphism or a subobject-classifier comparison is needed here.

For terminal objects, \(\partial1=T\) and \(d_1=1_T\).
The object \(A_1\) has the distinguished true member; on it
\(\widehat e_1(p,q)=\widehat d_1(p)\wedge\widehat d_1(q)=1\).
Its quotient is therefore \(1_0\). For products, form the paired
partial member \(\mu(p,q)\) on the intersection of the two
domains. It has
\[
 d_{X\times Y}\mu(p,q)=d_X(p)\wedge d_Y(q),\qquad
 e_{X\times Y}(\mu(p,q),\mu(p',q'))
                         =e_X(p,p')\wedge e_Y(q,q').       \tag{2}
\]
These follow by the two membership pullbacks, not an external choice
of values. Hence \(u_*\mu\) induces
\(H_hX\times H_hY\to H_h(X\times Y)\), using the effective
epi cover \(A_X\times A_Y\). Its inverse is the pair of
\(H_h\)'s projections. Projecting a paired partial member merely
restricts its original member to the common domain, whose h-value
is true. The resulting common-member predicate therefore has h-value
true. Conversely pairing the two projections of a partial pair
recovers the same partial pair. Thus both composites are identity
after the epi covers, proving the actual inverse maps. Product
naturality, symmetry and associativity follow by the same membership
identities and epi cancellation.

For any mono \(m:Z\to X\), equality in (1) proves that
\(H_hm\) is monic: lift two allegedly equal classes along the
quotient epis, use the target kernel pair, then the source kernel
pair. Now take \(m\) to be the equalizer of \(f,g:X\to Y\).
Restrict a partial member \(p\) to that equalizer, obtaining an
actual partial member \(\rho_m(p)\) of \(Z\). Its domain is
\[
             d_Z\rho_m(p)=e_Y(\partial f(p),\partial g(p)).
                                                               \tag{3}
\]
Indeed the two image members have a common value exactly on the
domain of \(p\) where \(f(x)=g(x)\). Over the pullback in
\(A_X\) of the equalizer of \(H_hf,H_hg\), effectiveness says
that the h-value of (3) is true. Thus \(u_*\rho_m\) gives an
h-total lift into \(A_Z\). Its image under \(m\) agrees with
\(p\) on that full h-domain. Lifts with the same quotient value
remain related by \(R_Z\), since \(m\) is monic. Descent
therefore constructs an inverse to \(H_hZ\) onto the equalizer.
The quotient cover of that equalizer is epi by pullback stability,
so these two inverse identities are genuine map identities. This
proves finite-limit preservation, including generalized parameters;
no Boolean logic or choice was used.

For an arbitrary selected mono \(m:Z\to X\), set
\(\gamma_m=\forall_{X\to1}(m)\) as in the exact localization.
The partial restriction still exists and satisfies
\[
 \gamma_m\wedge d_X(p)\le d_Z\rho_m(p),\qquad
 \gamma_m\wedge e_X(p,q)\le e_Z(\rho_m(p),\rho_m(q)).    \tag{4}
\]
On \(\gamma_m\), every member of \(X\) belongs to \(Z\);
this is the counit inequality of the already proved universal
quantifier, not commutation of \(h\) or \(u_*\) with it.
Selection means precisely that \(h\,u_*\chi_{\gamma_m}=1\).
Equations (4) then show that restriction gives \(A_X\to A_Z\)
and respects the equivalence relations. Its quotient map is inverse
to \(H_hm\): the common-value domain of \(m\rho_m(p)\)
and \(p\) has h-value true, and the other composite follows by
monicity. Thus every original source denominator is inverted.

The ordinary localization universal property of
Theorem~\ref{thm:p10-lawvere-localization-fractions} gives the displayed
\(\overline H_h\), with its exact fraction formula, equality on
refinements, functor laws and unique extension. To check finite limits
of genuinely localized maps, take a common admitted subterminal \(U\)
for a finite diagram and its finitely many equations. Its representatives
are a finite diagram over \(U\), and the localized limit is its
ordinary limit followed by the identification of \(Q(U\times X)\)
with \(QX\). We have \(H_hU\cong1_0\), since \(U\to1\)
is itself selected. Applying the proved finite-limit functor identifies
the transported diagram and its transported limit with the required
diagram in \(\mathcal E_0\). Hence the extension is left exact.
When the base is degenerate, all h-values are true and every quotient
is uniquely terminal; the original filter is improper and the ordinary
localized category is also degenerate. All constructions and equations
above remain valid in that boundary case.
\end{proof}

\annexitem{annex:p10.lawvere-closed-base.closed}{thm:p10-lawvere-closed-base-closed}{Evaluation, the closed axioms and the base hom category}
\begin{proof}
Put \(F=\overline H_h\). In the ordinary localization, the
exponential is the already proved \(Q(Y^X)\); it is not a newly
postulated hom object. The inverse product comparison and transported
evaluation from Theorem~\ref{thm:p10-lawvere-localization-logical} give
\[
 F(Y^X)\times FX\cong F(Y^X\times X)
                       \xrightarrow{F\operatorname{ev}}FY.
                                                               \tag{5}
\]
Define \(\kappa_{X,Y}\) as its unique exponential transpose in
\(\mathcal E_0\). This proves its actual existence and uniquely
fixes the comparison among product-induced comparisons. Naturality in
both arguments follows by transposing the natural evaluation equations;
the arguments may be arbitrary localized arrows, since \(F\) is
already a functor on those arrows. The unit map \(1_0\to F1\)
is the inverse of the proved terminal comparison.
Concretely, on the effective epi cover
\(A_{Y^X}\times A_X\), evaluation sends representatives \((p,x)\)
to the class of \(u_*\partial\operatorname{ev}(u_*\mu(p,x))\).
Equations (1)--(2) of the preceding proof ensure h-totality and
independence of both representatives. Thus (5) is an explicit
partial-member operation before its exponential transpose, not a
postulated interpretation of function values.

We check the historical closed-functor compatibilities, not just a
map named a comparison. The transported name \(Fj_X:F1\to F(X^X)\)
followed by \(\kappa\) evaluates on \(FX\) as its identity;
transposition gives the target identity name. This is CF1. For
\(i_X:X\cong X^1\), (5), restricted along the unit map, evaluates
as the identity of \(FX\). Its transpose is the canonical target
\(i_{FX}\), which is CF2. For CF3, untranspose its two routes
on
\(F(Z^Y)\times F(Y^X)\times FX\). Both are the transported
map which evaluates first \(f\) on \(x\), then \(g\) on that
value, because the source internal composition is characterized by
\(\operatorname{ev}(g\circ f,x)=\operatorname{ev}(g,
\operatorname{ev}(f,x))\). The product comparison is coherent by
(2) and functoriality. Thus the two maps agree; uniqueness of the
two transpositions gives precisely CF3. This proves an actual closed
functor in the finite selected 1966 definition. It does not assert its
comparison is invertible or its element functor normal.

The same calculation provides the base-enriched category explicitly.
Set \(\mathcal C_h(X,Y)=F(Y^X)\), take identity
\(1_0\cong F1\xrightarrow{Fj_X}F(X^X)\), and composition
\[
 F(Z^Y)\times F(Y^X)\cong F(Z^Y\times Y^X)
                                      \xrightarrow{Fc}F(Z^X).
\]
The source internal composition has associative and identity equations,
as may be checked by evaluation on \(g,f,x\); transport and product
coherence prove those same equations in the base. Equation (5) and CF3
make \(X\mapsto FX\), with hom map \(\kappa\), a base-enriched
functor into the self-enriched \(\mathcal E_0\). These are actual
hom objects, units, composition and comparisons at generalized bases.
On representatives the composite is the class of
\(u_*\partial c(u_*\mu(g,f))\), and the identity representative
is the transported singleton of the source identity name. Their
domains have h-value true by the same product-domain equation.

For clarity the ordinary naming maps are also explicit. A fraction
\([U,f]:X\to Y\) transposes to \(U\to Y^X\). As \(U\)
is subterminal, its graph is a partial member of \(Y^X\), with
domain \(U\). Its image under \(u_*\) is an h-total global
member of \(A_{Y^X}\), hence gives a global element of
\(H_h(Y^X)=\mathcal C_h(X,Y)\). Two such names are equal
exactly when the h-value of their common-member domain \(V\)
is true, by effectiveness. This domain is the intersection of their
two denominators on which their two transposed maps agree. Agreement
of the transposes is agreement of their evaluations on every \(x\),
by the established exponential bijection. Thus \(V\) is admitted
exactly when their fraction classes agree: one direction uses an
admitted common restriction below \(V\), and the other takes
\(V\) itself. Consequently the naming maps are injective and
preserve identities and composition. They are not claimed surjective.

The coequalizer property of \(A_X/R_X\) fixes its maps, and the
localization property fixes \(F\) once \(H_h\) is constructed.
The exponential transpose property fixes (5)'s comparison and supplies
the stated closed universal property. None of these is an adjoint-functor
assertion or a universal claim among all arbitrarily equipped closed
functors. In particular no normal-element-functor identity is smuggled
into these universal properties.
\end{proof}

\annexitem{annex:p10.lawvere-closed-base.countercases}{lem:p10-lawvere-closed-base-countercases}{Descent and enriched global elements are different requirements}
\begin{proof}
In sets, a subsingleton predicate on \(X\) is either empty or
\(\{x\}\), so \(\partial X=1+X\). For
\(\mathcal E=\mathsf{Set}\times\mathsf{Set}\), its partial
object at \((X,Y)\) is \((1+X,1+Y)\). Product \(u_*\)
sends this to \((1+X)\times(1+Y)\). First-coordinate \(h\)
makes h-total members exactly \(X\times(1+Y)\); the relation
identifies precisely those with the same first coordinate. Its quotient
is \(X\), even when \(Y\) is empty, since \(1+Y\) still
contains the empty partial member. Forward images act on the first
coordinate, so \(H_h\) is exactly first projection. The admitted
subterminal \((1,0)\to(1,1)\) is sent by product \(u_*\)
to \(0\to1\), but by \(H_h\) to an isomorphism. This proves
both the obstruction to original descent and the concrete new repair.

For diagonal \(u_*\), the source truth object maps to \((2,2)\),
already the target truth object, so take \(h\) identity. The only
admitted global subterminal of sets is \(1\); localization does
not change sets. In each component h-total partial members are exactly
singletons, and their equivalence is equality. Thus \(H_hX=(X,X)\)
on objects and maps. The closed comparison is componentwise the ordinary
set-exponential isomorphism. Nevertheless a global member of
\((Y^X,Y^X)\) is an arbitrary pair of functions, whereas an ordinary
map's transported name is \((f,f)\). For \(X=1,Y=2\), there
are four such global members but only two ordinary maps. The naming
injection is not surjective, and the ordinary category cannot be
identified with that enriched global underlying category. In particular
\(\Gamma_{\mathcal E_0}H_hX=X\times X\) does not equal
\(\Gamma_{\mathcal E_h}X=X\) generally. The separate normality
condition is genuinely stronger, not an automatic meaning of the closed
comparison or of the original printed data.
\end{proof}

\annexitem{annex:p10.lawvere-closed-base.collapse}{thm:p10-lawvere-closed-base-collapse}{A bounded Boolean ultrafilter and the two-valued target}
\begin{proof}
Theorem~\ref{thm:p10-lawvere-bvm-presentation} establishes \(B=p_*T\)
as an internally complete Boolean algebra in \(\mathcal S\), and
Lemma~\ref{lem:p10-lawvere-bvm-base} supplies bounded well-ordering
in this base. Nondegeneracy of \(\mathcal E\)
makes the two constants of \(B\) distinct: by adjunction their
equality would be equality of the two characteristic maps in
\(\mathcal E\). In \(\mathcal S\), the equalizer of these
constants is a subterminal, hence either zero or one; it is not one,
so \(B\) is internally nontrivial. Choose an internal well-order
of its underlying object.

We construct a proper deciding filter with bounded objects, rather
than assume a prime-ideal or ultrafilter theorem. Proper filters are
a definable subobject of \(\mathcal P_{\mathcal S}B\).
For a proper filter \(F\) and \(b\in B\), the generated
extension by \(b\) is explicitly
\[
       F[b]=\{z:\exists a\in F, a\wedge b\le z\}.
\]
It is a filter: upward closure is immediate, \(1\) belongs, and
two witnesses can be replaced by their meet in \(F\).
It is improper exactly when some \(a\in F\) has \(a\wedge b=0\).
At least one of \(F[b],F[\neg b]\) is proper, since witnesses
for impropriety of both would give \(a\le\neg b\),
\(a'\le b\), and \(a\wedge a'=0\in F\).
By classicality, prescribe \(F[b]\) if it is proper and otherwise
prescribe \(F[\neg b]\). This is a specified operation, not a
new choice principle.

Here is the existence proof of the resulting well-ordered history.
A history on an initial segment \(D\subseteq B\) is a map
\(b\mapsto F_b\) to proper filters, where \(F_b\) is that
prescribed extension of
\(\{1\}\cup\bigcup_{a<b}F_a\). This union is a proper
filter: the preceding filters form an inclusion chain; for a binary
meet take the larger of two containing filters, and zero belongs
to none. The monotonicity follows from the prescribed extension.
The empty history is allowed. Two histories agree on their common
domain, since a least disagreement would have identical earlier
filters and hence the same prescribed extension. Their domains are
nested initial segments.

The collection of all such domain/graph pairs is a subobject of
\(\mathcal P(B)\times\mathcal P(B\times\mathcal P(B))\).
Its union is therefore an actual functional graph on an initial
segment \(D\), and is itself a history: for any member \(b\),
a constituent containing it contains all its predecessors, with
the same filters. If \(D\ne B\), take the least missing \(b\).
The union of its preceding filters is again proper by the chain
argument; its prescribed extension appends a new member, producing
a history outside the supposedly total union. This is a contradiction.
Thus the history has domain \(B\). No external ordinal recursion,
Replacement, infinite product, or unproved recursion-lifting theorem
was used.

Set \(\mathcal U=\{1\}\cup\bigcup_bF_b\). It is a proper
filter and decides each \(b\): exactly one of \(b,\neg b\)
belongs. The classifier \(h:B\to2=T_{\mathcal S}\) of
\(\mathcal U\) preserves \(0,1\), meets and complements.
For joins, if neither \(b\) nor \(c\) belongs, their complements
belong, so \(\neg(b\vee c)\) belongs; conversely a member
\(b\) or \(c\) makes the join a member. Thus it preserves
binary joins, and it preserves implication because Boolean implication
is \(\neg b\vee c\). It is exactly a Heyting homomorphism of
the required type. Arbitrary-join preservation is neither proved nor
required; an atomless complete Boolean algebra can have no such
complete homomorphism to \(2\). Indeed for a complete homomorphism
the meet \(a\) of its true filter would still have value true,
so \(a\ne0\). For \(c\le a\), either \(c\) has value true,
forcing \(a\le c\), or its complement has value true, forcing
\(a\le\neg c\) and \(c=0\). Thus \(a\) would be an atom.

Apply the full original datum \((p_*,h)\). Its global truth filter
is proper since \(h(0)=0\ne1\). Hence the established logical
localization is nondegenerate: if zero became one, its eventual
characteristic equality would put zero in this proper filter.
Every subterminal of \(\mathcal E_h\) is represented, after an
admitted restriction, by a strict subterminal of \(\mathcal E\),
by the proved eventual-mono result. For each such \(V\),
\(h(p_*\chi_V):1\to2\) is either true or false, since
\(\mathcal S\) has only two subterminals. In the first case
\(QV=1\). In the second, its Boolean complement is admitted,
so \(QV=0\), using logical preservation of that complementary
sum. Thus precisely zero and one are subterminals in the nondegenerate
target. Classicality follows from \(QT\cong Q(1+1)\cong1+1\).

For completeness choice also survives the genuinely localized maps.
Represent an epi by a numerator over an admitted \(U\).
Its image inclusion becomes an isomorphism in the localization.
The universal truth of that inclusion is therefore admitted: equality
of its classifier with true on a further admitted restriction is
exactly this condition, by the fraction equality criterion and the
internal universal-quantifier bijection. Restrict to the intersection
with that admitted truth. The numerator over this smaller subterminal
is now an ordinary epi in the corresponding slice of \(\mathcal E\).
It splits there: an underlying epi has a section by
Theorem~\ref{thm:p10-lawvere-bvm-realization}, and its section lies over the same
base since its composite with the epi is identity. Localizing this
section gives a section of the original germ epi. Thus every epi
in \(\mathcal E_h\) splits. Together with classicality this
retains the source Boolean package; subterminal generation also makes
the two-valued target well pointed. The constructed closed functor
to \(\mathcal S\) needs no geometric adjoint assertion.

In this special case we can also prove, rather than assume, normality
of the element comparison. A global element of \(H_hZ\) lifts
along \(A_Z\to H_hZ\), since every epi in \(\mathcal S\)
splits. The lift is a global element of \(p_*\partial Z\),
which adjunction identifies with a global partial member of \(Z\)
in \(\mathcal E\), because \(p^*1=1\). Its domain \(U\)
is admitted by h-totality, so that partial member is exactly a
fraction point \([U,z]:1\to Z\). The injective naming argument
therefore proves a natural bijection
\(\mathcal E_h(1,Z)\cong\mathcal S(1,H_hZ)\).
Applied to every exponential this also identifies ordinary hom sets
with the global elements of \(\mathcal C_h(X,Y)\). Identities
and composition agree by their already proved naming equations.
Moreover (5) takes the name of a fraction \(f\) through \(\kappa\)
to the name of the actual map \(Ff\), by evaluation and transposition.
Thus the hom comparison satisfies the normality compatibility NCF1
as well as the element-functor identification NCF0.
Transporting the element functors through these canonical bijections
gives the separate normality requirement. We have used the actual
geometric adjunction and epi splitting of this particular base;
neither is present for an arbitrary original \(u_*\), so the
diagonal countercase is not contradicted.

Finally if the original model is degenerate, \(B\) is the
one-element algebra. A homomorphism preserving its coincident zero
and one cannot land in the distinct zero and one of \(2\).
This supplies the necessary exception to an unqualified reading of
``always'' in the original collapse sentence. The full construction
with a degenerate base instead gives the separately retained improper
boundary, not a nondegenerate two-valued collapse.
\end{proof}
